\documentclass[10pt]{article}
\usepackage[T1]{fontenc}
\usepackage[utf8]{inputenc}
\usepackage[french]{babel}
\usepackage{amsmath,amssymb,mathrsfs,mathpazo}
\usepackage{geometry}
\geometry{a4paper,margin=2cm}
\usepackage{graphicx,array,tabularx,tabularray,colortbl,multirow,multicol,booktabs,enumitem,pifont,ulem,fancybox,diagbox,cancel}
\usepackage{pstricks,pstricks-add,pst-plot,pst-tree,pst-node,pst-text,pst-func,pst-math,pst-bspline,pst-3dplot,pst-eucl,pst-solides3d}
\let\pscalculate\undefined
\usepackage{pst-bezier}
\usepackage{tikz,pgfplots,tkz-tab}
\usetikzlibrary{arrows,arrows.meta,patterns,calc,positioning,trees,decorations.markings,shapes.geometric,shapes.misc,shapes.arrows,intersections,decorations.pathmorphing,decorations.pathreplacing,backgrounds,shadows,babel}
\pgfplotsset{compat=1.18}
\usepgfplotslibrary{fillbetween}
\usepackage[np]{numprint}
\usepackage{listings,makecell,pas-tableur,twoopt}
\let\barre\undefined
\usepackage[pstricks]{tablvar}
\definecolor{darkgreen}{RGB}{0,100,0}
\providecommand{\red}{\color{red}}
\providecommand{\blue}{\color{blue}}
\providecommand{\green}{\color{green}}
\providecommand{\black}{\color{black}}
\providecommand{\white}{\color{white}}
\providecommand{\gray}{\color{gray}}
\providecommand{\cyan}{\color{cyan}}
\providecommand{\magenta}{\color{magenta}}
\providecommand{\yellow}{\color{yellow}}
\providecommand{\violet}{\color{violet}}
\providecommand{\orange}{\color{orange}}
\providecommand{\euro}{\texteuro}
\providecommand{\EUR}{\texteuro}
\providecommand{\ieme}{\textsuperscript{e}}
\providecommand{\degres}{\ensuremath{^\circ}}
\providecommand{\degree}{\ensuremath{^\circ}}
\providecommand{\vv}[1]{\overrightarrow{#1}}
\providecommand{\Coord}[1]{#1}
\providecommand{\decosix}{}
\providecommand{\starredbullet}{\ensuremath{\bullet}}
\pagestyle{empty}
\setlength{\parindent}{0pt}
\usepackage{adjustbox}\DeclareUnicodeCharacter{2236}{:}\DeclareUnicodeCharacter{2190}{\ensuremath{\leftarrow}}
\providecommand{\R}{\mathbb{R}}
\providecommand{\N}{\mathbb{N}}
\providecommand{\Z}{\mathbb{Z}}
\providecommand{\Q}{\mathbb{Q}}
\providecommand{\D}{\mathbb{D}}
\providecommand{\C}{\mathbb{C}}
\providecommand{\e}{{\mathrm{e}}}
\providecommand{\vect}[1]{\overrightarrow{#1}}
\providecommand{\vectt}[1]{\overrightarrow{\mathrm{#1}}}
\providecommand{\vvt}[1]{\overrightarrow{\mathrm{#1}}}
\providecommand{\np}[2][]{{#2~#1}}
\providecommand{\vv}[1]{{\overrightarrow{#1}}}
\renewcommand{\d}{{\mathrm{d}}}
\providecommand{\degres}{{^\circ}}
\providecommand{\degre}{{^\circ}}
\providecommand{\blue}{\color{blue}}
\providecommand{\red}{\color{red}}
\providecommand{\enskip}{\quad}
\providecommand{\llbracket}{\lbrack\!\lbrack}
\providecommand{\rrbracket}{\rbrack\!\rbrack}
\providecommand{\euro}{{\text{€}}}
\providecommand{\no}{{\text{n°}}}
\providecommand{\interval}[2]{{[#1;#2]}}
\providecommand{\intervFO}[2]{{[#1;#2[}}
\providecommand{\dotfill}{\ldots}

\providecommand{\ds}{\displaystyle}
\providecommand{\textding}[1]{#1}
\providecommand{\Oij}{\(\mathrm{O};\vec{\imath},\vec{\jmath}\)}
\providecommand{\Oijk}{\(\mathrm{O};\vec{\imath},\vec{\jmath},\vec{k}\)}
\providecommand{\Ouv}{\(\mathrm{O};\vec{u},\vec{v}\)}
\providecommand{\intervFF}[2]{\ensuremath{[#1;#2]}}
\renewcommand{\euro}{\ifmmode\text{\texteuro}\else\texteuro\fi}
\DeclareUnicodeCharacter{00B0}{\ifmmode^{\circ}\else\textdegree\fi}
\DeclareUnicodeCharacter{20AC}{\ifmmode\text{\texteuro}\else\texteuro\fi}
% Énoncés : Éducation nationale. Transcription et corrigés : APMEP, Denis Vergès et contributeurs. Découpage et classement : CoopMaths, Sébastien Lozano, Éric Elter et contributeurs. Adaptation de la présentation : Maths Explorer.
% CC BY-SA : https://coopmaths.fr/www/about/
% https://coopmaths.fr/alea/static/eam/2026/tex/eam_technologique_2026_09_antilles_automatismes.tex
\begin{document}
\noindent\textit{APMEP / CoopMaths. CC BY-SA. Présentation adaptée par Maths Explorer.}\par\medskip
\textbf{Pour cette première partie, aucune justification n'est demandée et une seule réponse est possible par question. Pour chaque question, reportez son numéro sur votre copie et indiquez votre réponse.}

\needspace{10\baselineskip}
\subsection*{Question 1}

Un pantalon coûtant initialement 80~\euro{} est soldé à 25\,\%. Quel est le nouveau prix ?

\medskip
\begin{tblr}{colspec = {@{}*{4}{X}}}
	\textbf{A.}\quad 20~\euro &
	\textbf{B.}\quad 25~\euro &
	\textbf{C.}\quad 55~\euro &
	\textbf{D.}\quad 60~\euro \\
\end{tblr}

\needspace{10\baselineskip}
\subsection*{Question 2}

Paul dépense 600~\euro{} dans la location de son appartement, ce qui représente 30\,\% de son salaire. Le salaire de Paul est de :

\medskip
\begin{tblr}{colspec = {@{}*{4}{X}}}
	\textbf{A.}\quad 180~\euro &
	\textbf{B.}\quad \np{1800}~\euro &
	\textbf{C.}\quad \np{2000}~\euro &
	\textbf{D.}\quad \np{4200}~\euro \\
\end{tblr}

\needspace{10\baselineskip}
\subsection*{Question 3}

La fréquentation d'une salle de cinéma au mois de juillet a diminué de 8\,\% par rapport au mois de juin. On a compté 532 spectateurs en juillet, le nombre de spectateurs en juin était donc de :

\medskip
\begin{tblr}{colspec = {@{}*{4}{X}}}
	\textbf{A.}\quad $532 \times 0,92$ &
	\textbf{B.}\quad $532 \div 0,92$ &
	\textbf{C.}\quad $532 \times 1,08$ &
	\textbf{D.}\quad $532 \div 1,08$ \\
\end{tblr}

\needspace{10\baselineskip}
\subsection*{Question 4}

Quelle est la plus petite de ces quatre fractions : $\dfrac{1}{4}$, $\dfrac{3}{4}$, $\dfrac{1}{8}$, $\dfrac{2}{15}$ ?

\medskip
\begin{tblr}{colspec = {@{}*{4}{X}}}
	\textbf{A.}\quad $\dfrac{1}{4}$ &
	\textbf{B.}\quad $\dfrac{3}{4}$ &
	\textbf{C.}\quad $\dfrac{1}{8}$ &
	\textbf{D.}\quad $\dfrac{2}{15}$ \\
\end{tblr}

\needspace{10\baselineskip}
\subsection*{Question 5}

Parmi les nombres suivants, lequel est différent de 10,25 ?

\medskip
\begin{tblr}{colspec = {@{}*{4}{X}}}
	\textbf{A.}\quad $\np{0,1025} \times 10^{-2}$ &
	\textbf{B.}\quad $\dfrac{\np{1025}}{100}$ &
	\textbf{C.}\quad $\dfrac{205}{20}$ &
	\textbf{D.}\quad $102,5 \times 10^{-1}$ \\
\end{tblr}

\needspace{10\baselineskip}
\subsection*{Question 6}

\begin{center}
	\textbf{\footnotesize\scshape Graphique. Liens entre la pratique et l'âge}\\
	{\footnotesize\scshape L'exemple des 8 activités issues des univers les plus prisés}

	\medskip

	\fbox{\sffamily\small\begin{tikzpicture}[y=0.09cm,x=1.15cm]
		\def\cats{Football/16/3,Fitness/13/5,{Footing,\\jogging}/25/11,Tennis/8/4,VTT/11/6,{Nage,\\natation}/22/15,{Vélo sur\\route}/13/8,{Randonnée\\pédestre}/17/27}
		\foreach [count=\i from 0] \nom/\jeune/\age in \cats{
			\fill[teal] (\i,0) rectangle++ (0.35,\jeune);
			\node[font=\tiny,above,inner sep=2pt] at (\i+0.175,\jeune) {\jeune\,\%};
			\fill[orange] (\i+0.4,0) rectangle++ (0.35,\age);
			\node[font=\tiny,above] at (\i+0.575,\age) {\age\,\%};
			\node[font=\tiny,align=center,below] at (\i+0.4,-0.5) {\nom};
		}
		\node[font=\tiny] at (3.9,-10) {\textcolor{teal}{$\blacksquare$}~Moins de 40 ans \qquad \textcolor{orange}{$\blacksquare$}~40 ans et plus};
	\end{tikzpicture}}

	{\tiny\sffamily\textbf{Source} : INJEP, Ministère des sports, CRÉDOC : Baromètre national des pratiques sportives 2018}\\
	{\tiny\sffamily\textbf{Champ} : ensemble de la population de 15 ans et plus résidant en France DROM compris}
\end{center}

D'après le graphique ci-dessus on peut affirmer que, parmi les sportifs de moins de 40 ans, 11\,\% pratiquent :

\medskip
\begin{tblr}{colspec = {@{}*{4}{X}}}
	\textbf{A.}\quad le VTT &
	\textbf{B.}\quad le footing-jogging &
	\textbf{C.}\quad le vélo sur route &
	\textbf{D.}\quad la randonnée pédestre \\
\end{tblr}

\needspace{10\baselineskip}
\subsection*{Question 7}

Une expression d'une fonction affine ayant le tableau de signes ci-dessous

\begin{center}
	\begin{tikzpicture}
		\tkzTabInit[espcl=2.5]{$x$/0.8,signe de $f(x)$/1}{$-\infty$,$3$,$+\infty$}
		\tkzTabLine{,-,z,+,}
	\end{tikzpicture}
\end{center}

est :

\medskip
\begin{tblr}{colspec = {@{}*{4}{X}}}
	\textbf{A.}\quad $f(x) = -3x+9$ &
	\textbf{B.}\quad $f(x) = 2x-6$ &
	\textbf{C.}\quad $f(x) = -x-3$ &
	\textbf{D.}\quad $f(x) = 3x-10$ \\
\end{tblr}

\needspace{10\baselineskip}
\subsection*{Question 8}

\begin{minipage}[t]{0.52\linewidth}
	On donne ci-contre la courbe représentative d'une fonction $f$ définie sur l'intervalle $[-4 ; 6]$. Choisir la bonne réponse.

	\medskip
	\begin{tblr}{colspec = {@{}X}}
		\textbf{A.}\quad $f$ est positive sur l'intervalle $[4 ; 6]$ \\
		\textbf{B.}\quad $f$ est croissante sur l'intervalle $[-5 ; 3,5]$ \\
		\textbf{C.}\quad $f$ est positive sur l'intervalle $[2 ; 5]$ \\
		\textbf{D.}\quad $f$ est croissante sur l'intervalle $[-4 ; -1]$ \\
	\end{tblr}
\end{minipage}
\hfill
\begin{tikzpicture}[x=0.5cm,y=0.5cm,>=stealth,baseline={(current bounding box.north)}]
	\draw[step=1,gray!75,line width=0.5pt] (-5,-5.5) grid (6.5,4.5);
	\draw[->,line width=1pt] (-5,0) -- (6.5,0);
	\draw[->,line width=1pt] (0,-5.5) -- (0,4.5);
	\foreach \x in {-5,-4,-3,-2,-1,1,2,3,4,5,6} \draw (\x,2pt) -- (\x,-2pt) node[below] {\scriptsize \np{\x}};
	\foreach \y in {-5,-4,-3,-2,-1,1,2,3,4} \draw (2pt,\y) -- (-2pt,\y) node[left] {\scriptsize \np{\y}};
	\draw[line width=1.2pt,red] plot[domain=-4:6,samples=100] (\x,{0.1*(\x+3)*(\x-2)*(\x-5)});
\end{tikzpicture}

\needspace{10\baselineskip}
\subsection*{Question 9}

La formule $E = mc^2$ établie par Albert Einstein en 1905, relie la vitesse de la lumière dans le vide $c$, l'énergie $E$ et la masse non nulle $m$ d'une particule.

On peut aussi écrire que :

\medskip
\begin{tblr}{colspec = {@{}*{4}{X}}}
	\textbf{A.}\quad $c = \dfrac{E}{m} \div 2$ &
	\textbf{B.}\quad $m = E - c^2$ &
	\textbf{C.}\quad $c = \sqrt{E-m}$ &
	\textbf{D.}\quad $c = \sqrt{\dfrac{E}{m}}$ \\
\end{tblr}

\needspace{10\baselineskip}
\subsection*{Question 10}

En vous aidant d'un ordre de grandeur, choisir la bonne réponse.

$289 \times 0,307 = $

\medskip
\begin{tblr}{colspec = {@{}*{4}{X}}}
	\textbf{A.}\quad \np{0,88723} &
	\textbf{B.}\quad \np{8,8723} &
	\textbf{C.}\quad \np{88,723} &
	\textbf{D.}\quad \np{887,23} \\
\end{tblr}

\needspace{10\baselineskip}
\subsection*{Question 11}

$(2x-7)^2$ est égale à :

\medskip
\begin{tblr}{colspec = {@{}*{4}{X}}}
	\textbf{A.}\quad $4x^2-28x+49$ &
	\textbf{B.}\quad $4x^2+28x-49$ &
	\textbf{C.}\quad $2x^2-28x+49$ &
	\textbf{D.}\quad $2x^2+28x-49$ \\
\end{tblr}

\needspace{10\baselineskip}
\subsection*{Question 12}

La durée correspondant à 1,2~h est :

\medskip
\begin{tblr}{colspec = {@{}*{4}{X}}}
	\textbf{A.}\quad 62 minutes &
	\textbf{B.}\quad 72 minutes &
	\textbf{C.}\quad 80 minutes &
	\textbf{D.}\quad 120 minutes \\
\end{tblr}

\end{document}
